\section{Setup}

All models share one recipe from the course codebase: vocabulary 4096, context 1024, width 512, batch 64 sequences, 600{,}000 training sequences (614.4M tokens, 9{,}374 steps), AdamW with weight decay 0.1, 1\% linear warmup followed by linear decay to zero, gradient clipping at norm 1, mixed precision training, model seed 42 and data seed 42. The metric throughout is the validation loss on the held out DCLM split after the last step (lower is better). Peak learning rate is the only hyperparameter swept. P1 and P2 use the 105 supplied reference runs (21 variants, five rates each, centred on the instructor's trial rate). P3 and P4 runs were trained for this report on H100s in my Modal environment with the same recipe and are listed with their rates in each section.

\paragraph{Diagnostics.} Checkpoint probes run the FP32 copy of each model on the same eight held out validation sequences (full backpropagation through time for the RNNs). For every branch input along the depth stream and for the final stream before the output norm I record the activation RMS, the gradient RMS of the full next token loss and of the loss on the final token only, the RMS ratio of the raw branch output to its input, and the cosine between consecutive stream states. For attention layers I record, for 32 query positions, the softmax mass on the first token, the mass on the query itself and the mean attention lag, which I compare with a uniform causal prefix baseline (mass $1/(t+1)$ on each of the $t+1$ visible keys, mean lag $t/2$). For recurrent layers I record per time step the hidden state RMS, the hidden state gradient RMS under both losses, the fraction of $|h|>0.95$ and the cosine between $h_t$ and $h_{t-\ell}$ for lags 1, 4, 16, 64. Context tests rerun the final token prediction with only the $k$ most recent input tokens (original positions preserved, recurrent state reset), with and without the first token restored, on 64 validation sequences. Three checkpoints are probed per selected run: the seed 42 initialisation, step 100 and the final model. All probe code is in \texttt{experiments/a4\_stability/measure\_all.py} and the plotting code in \texttt{report/plot\_*.py}; the figures only use these probe files and the W\&B logs.

\paragraph{Compute.} Training used about \budgethours{} of the 48 GPU hour budget. The reference P1/P2 probes ran on CPU containers because both GPU slots were busy with P3/P4 training; results are identical in FP32 up to accumulation order.
